\section{Related Work and Conclusion}\label{sec:conclusion}
There is a large amount of literature on change-point detection, see for example \cite{basseville1993detection,brodsky1993nonparametric}.
Optimal segmentation can be found using dynamic programming \cite{lavielle2006changepoint}; however, the complexity of this approach is quadratic in the number of samples $n$, and therefore might be infeasible for large data sets. Some approaches which achieve complexity linear in $n$ \cite{levy2007catching,killick2012optimal} treat only one dimensional data. %\ik{add refs to Bayesian methods}
Some related work is concerned with the objective \eqref{eq:segmentationObjective_mu} we presented in \secref{sec:algorithms}. In \cite{levy2007catching} it was suggested to reformulate \eqref{eq:segmentationObjective_mu} for the one dimensional case as a LASSO regression problem \cite{tibshirani1996regression,yuan2006model}, while %\cite{vert2010fast}
\cite{bleakley2011group}
extended this approach to multidimensional data, although not treating outliers directly.
%Our approach to optimizing \eqref{eq:segmentationObjective_mu} led to a different approximate solution.
%--- IK: uncomment the following sentence, if we put analysis back
%Furthermore, our analysis is more general as we do not assume a Gaussian data model as assumed in the analysis by \cite{vert2010fast}.
Another common approach is deriving an objective from a maximum likelihood criterion of a generative model, and then either optimize the objective or use it as a criterion for a top-down or a bottom-up approach \cite{qiao:2012study,phonemeSeg:2008,gracia:2011hierarchical,olshen2004circular}. We note that the two-dimensional version of \eqref{eq:segmentationObjective_mu} is used in image denoising applications, where it is known as the Total-Variation of the image \cite{Rudin:1992,Chambolle:2004,Beck_TV:2009}.
%As a denoising method it is known to preserve edges in the image, and indeed segmentation can be seen as an aggressive variant of denoising, where all data in a signal except for segment boundaries is considered as `noise' and is zeroed out.
%in the solution for \eqref{eq:segmentationObjective_mu}.
Finally, we note that all these approaches do not directly incorporate outliers into the model.

%The following related works are focused on convex clustering, but are nevertheless closely related to the convex segmentation task. Lashkari and Golland \cite{Lashkari07convexclustering} formulated the convex clustering using likelihood function for a mixture model, where the mixture components are centered around the data samples. Lindsten et al.~\cite{lindsten:2011a, lindsten:2011b} used a couple of relaxation steps of k-means clustering where $l_0$ norm was replaced by $l_1$ norm, to derive a convex objective with sum of $l_p$ norms regularization. Convex relaxation of hierarchical clustering and optimization schemes for finding the solution paths was presented by Hocking et al.~\cite{hocking:2011clusterpath}, for $l_p$ regularization with $p=1,2,\infty$. Optimization methods in the specific context of convex clustering was presented by Chi and Lange \cite{Chi:2013arXiv}.

We formulated the task of segmenting sequential data and detecting outliers using convex optimization, which can be solved in an alternating manner. We showed that a specific choice of weighting
%recovers the solution to an unrelaxed optimization problem for the case of $K=2$ segments, and
can empirically enhance performance.
We described how to calculate
$\lambda^*$ and $\gamma^*$, the critical values for the split into two
segments and the detection of the first outlier, respectively. These
values are useful for finding a solution path in the two-dimensional
parameter space. We also derived a top-down, outlier-robust
hierarchical segmentation algorithm which minimizes the objective in a
greedy manner.
%We proposed an optimization approach for minimizing the objective, derived an analytical solution for the case of $K=2$ segments, and used this solution to derive a hierarchical top-down segmentation algorithm.
%
%We showed that under the assumption that two sub-sequences are generated from two different processes, the segments boundary detected by the algorithm is not too far from the correct boundary. We used the derived bound in order to incorporate weighting into the algorithm, which result in improved performance.
%We further proposed an extension which improves
%performance under non-stationary noise, by incorporating outlier-robustness
%into the model.
%The resulting optimization problem is convex as well and can
%be solved in an alternating manner.
This algorithm allows for directly controlling both the
number of desired segments $K$ and number of outliers $M$. Experiments with %synthetic data and
real-world audio data with outliers added manually demonstrated the superiority of our algorithms.

We consider a few possible extensions to the current work. One is deriving algorithms that will work on-the-fly. Another direction is to investigate more involved noise models, such as noise which corrupts a single feature along all samples, or a consecutive set of samples. Yet another interesting question is how to identify that different segments come from the same source, e.g.\ that the same speaker is present at different locations in a recording. We plan to investigate these directions in future work.
%
%\begin{itemize}
%    \item Treat more involved noise models, such as noise that corrupt a single feature along all samples, or noise that corrupt a consecutive set of samples.
%    \item Derive algorithms that will work on-the-fly.
%    \item Add the ability to detect different segments as originating from the same source, for example identify a specific speaker in two different segments.
%    \item Learn the structure of samples within a segment, by considering different distance measures other than Euclidean in the objective, as was described in \secref{sec:intra-segment-metric}.
%    \item Find the optimal values of the parameters $\lambda,\gamma$, or equivalently the correct number of segments and outliers, from the data.
%    \item Consider a kernel-formulation of the objectives.
%    \item Impose a structure on the sparsity pattern of the outliers variable, for example to get contiguous nonzero patterns.
%\end{itemize}
